\documentclass[11pt,letterpaper]{article}

\usepackage[utf8]{inputenc}
\usepackage[T1]{fontenc}
\usepackage[spanish,es-noshorthands]{babel}
\usepackage{spectral}
\usepackage[margin=2.5cm,headheight=48pt,headsep=18pt]{geometry}
\usepackage{fancyhdr}
\usepackage{graphicx}
\usepackage{xcolor}
\usepackage{listings}
\usepackage{tikz}
\usepackage{booktabs}
\usepackage{enumitem}
\usepackage{float}
\usepackage[hidelinks]{hyperref}
\usetikzlibrary{positioning}

\newcommand{\curso}{CC4303-1 Redes}
\newcommand{\tituloinforme}{Actividad: Fragmentación}

\renewcommand{\sectionmark}[1]{\markboth{#1}{}}
\setlength{\parskip}{0.6em}
\setlength{\parindent}{0pt}
\setlength{\emergencystretch}{3em}

% Cabecera y pie de página con el mismo estilo del informe de práctica.
\pagestyle{fancy}
\fancyhf{}
\fancyhead[L]{\small\leftmark}
\fancyfoot[L]{\small Informe \tituloinforme}
\fancyfoot[R]{\small CC4303-1}
\renewcommand{\headrulewidth}{0.4pt}
\renewcommand{\footrulewidth}{0.4pt}

\fancypagestyle{portada}{%
  \fancyhf{}
  \fancyhead[L]{\small\begin{tabular}[b]{@{}l@{}}
    Universidad de Chile\\
    Facultad de Ciencias Físicas y Matemáticas\\
    Departamento de Ciencias de la Computación\\
    \curso
  \end{tabular}}
  \fancyhead[R]{\includegraphics[height=1.3cm]{logo_fcfm.png}}
  \renewcommand{\headrulewidth}{0.4pt}
  \renewcommand{\footrulewidth}{0pt}
}

\definecolor{fondocodigo}{gray}{0.96}
\lstset{
  basicstyle=\ttfamily\footnotesize,
  backgroundcolor=\color{fondocodigo},
  frame=single,
  rulecolor=\color{gray!50},
  breaklines=true,
  columns=fullflexible,
  keepspaces=true,
  showstringspaces=false,
  literate={á}{{\'a}}1 {é}{{\'e}}1 {í}{{\'i}}1 {ó}{{\'o}}1 {ú}{{\'u}}1
           {Á}{{\'A}}1 {ñ}{{\~n}}1
}

% Diagramas de red: routers como nodos y enlaces como líneas.
\tikzset{
  router/.style={circle,draw,thick,minimum size=8mm,font=\small},
  default/.style={rectangle,draw,thick,rounded corners,minimum size=8mm,font=\small,fill=gray!10},
  link/.style={thick,<->,>=stealth},
}

\begin{document}

\begin{titlepage}
  \thispagestyle{portada}
  \centering
  \vspace*{5cm}
  {\LARGE Informe\par}
  \vspace{0.3cm}
  {\Large \tituloinforme\par}
  \vspace{0.4cm}
  {\large Juan Pablo Bustamante Flaño\par}
  \vfill
  Correo electrónico: juan.bustamante.f@ug.uchile.cl\\
  Curso: \curso{} --- Primavera 2026\\
  Profesora: Ivana Bachmann\\
  Módulo 4: Redes y Ruteo, semanas 11-12\\
  Fecha: 18/11/2026
  \vspace*{3cm}
\end{titlepage}

\section{Introducción}

Se agregó fragmentación al mini-Internet de la actividad de forwarding básico (versión con
TTL). Cada enlace tiene un MTU: si un paquete no cabe, el router lo divide en fragmentos, y
el router de destino los re-ensambla antes de imprimir el mensaje. El código está en
\texttt{actividades/Fragmentacion} y se probó con la red de 5 routers y los MTU del
enunciado (figura~\ref{fig:red}).

\begin{figure}[H]
\centering
\begin{tikzpicture}[node distance=2.2cm]
  \node[router] (r1) {R1};
  \node[router,right=of r1] (r2) {R2};
  \node[router,right=of r2] (r3) {R3};
  \node[router,below=1.4cm of r2] (r4) {R4};
  \node[router,below=1.4cm of r3] (r5) {R5};
  \draw[link] (r1) -- node[above,font=\footnotesize] {100} (r2);
  \draw[link] (r2) -- node[above,font=\footnotesize] {50} (r3);
  \draw[link] (r2) -- node[left,font=\footnotesize] {80} (r4);
  \draw[link] (r3) -- node[right,font=\footnotesize] {55} (r5);
  \draw[link] (r4) -- node[below,font=\footnotesize] {140} (r5);
\end{tikzpicture}
\caption{Red de 5 routers con el MTU de cada enlace (bytes).}
\label{fig:red}
\end{figure}

\section{Diseño}

\textbf{Header (20 bytes).} IP (4), puerto (2), TTL (1), ID (4), offset (4), tamaño (4) y
FLAG (1), seguidos del mensaje. El ID y el offset usan 4 bytes porque el tamaño también usa
4 bytes, así que el offset debe poder indicar cualquier posición del mensaje. El FLAG usa un
byte completo con valor 0 o 1, que es más simple de leer y escribir que un bit suelto.

\textbf{Rutas.} \texttt{check\_routes} lee el MTU de la sexta columna de la tabla y retorna
\texttt{(next\_hop, MTU)}. Las líneas de la tabla de R5 del enunciado no traen el prefijo
\texttt{/24}; en ese caso la red se interpreta como una sola IP (\texttt{/32}).

\textbf{Fragmentación.} Si el paquete completo mide más que el MTU, el mensaje se divide
en trozos de \texttt{MTU - 20} bytes. El offset de cada fragmento se calcula respecto al
datagrama original, y el FLAG es 1 en todos los trozos menos el último, que conserva el
FLAG del paquete fragmentado. Así, al fragmentar un fragmento intermedio, su último trozo
mantiene FLAG 1.

\textbf{Re-ensamblaje.} Un paquete sin fragmentar tiene offset 0 y FLAG 0; cualquier otro
caso es un fragmento. \texttt{reassemble\_IP\_packet} ordena los fragmentos por offset,
toma el largo total del fragmento con FLAG 0 y concatena los mensajes. Retorna
\texttt{None} si falta un trozo. El router guarda en un diccionario, por ID, los fragmentos
dirigidos a él, intenta re-ensamblar con cada fragmento que llega y borra el ID cuando lo
logra.

\textbf{Tests de las funciones.} El header vuelve igual después de \texttt{parse\_packet}
y \texttt{create\_packet}. Un paquete de 150 bytes se divide en 2 fragmentos con MTU 100 y
en 5 con MTU 50, y \texttt{reassemble\_IP\_packet} lo recupera completo, incluso si los
fragmentos vienen desordenados o fueron re-fragmentados. Al enviar a mano a R5 los
fragmentos de dos paquetes con distinto ID, en orden, en desorden e intercalados, ambos
mensajes se re-ensamblaron correctamente.

\section{Pruebas}

Los paquetes se enviaron con \texttt{enviar.py}, que recibe todos los campos del header. No
se usó \texttt{netcat}, porque el header en bytes no se puede escribir como texto:

\begin{lstlisting}
% python3 enviar.py "127.0.0.1;8885;10;347;0;5;0" "hola!" 127.0.0.1 8881
\end{lstlisting}

\subsection{TTL y round-robin}

El paquete \texttt{hola!} se envió 4 veces de R1 a R5. Los 4 llegaron, el TTL bajó en 1 en
cada salto, y R2 siguió alternando entre R3 y R4:

\begin{lstlisting}
redirigiendo paquete 127.0.0.1;8885;8;347;0;5;0;hola! ... hacia ('127.0.0.1', 8883)
redirigiendo paquete 127.0.0.1;8885;6;347;0;5;0;hola! ... hacia ('127.0.0.1', 8884)
\end{lstlisting}

Los fragmentos también respetan el TTL: un paquete de 150 bytes enviado con TTL 2 se
fragmentó en R1 y en R2, y R3 y R4 descartaron todos sus fragmentos al llegar con TTL 0.

\subsection{Paquetes de 150 bytes}

Se envió 4 veces un paquete de 150 bytes de R1 a R5, y luego desde otros routers. Este es
el recorrido de uno de los envíos de R1 a R5:

\begin{lstlisting}
# R1 (enlace R1-R2 con MTU 100)
  fragmentado en 2 fragmentos para el enlace con MTU 100:
    127.0.0.1;8885;19;501;0;80;1;...
    127.0.0.1;8885;19;501;80;50;0;...
# R2: el fragmento de 100 bytes va hacia R3 (MTU 50) y se vuelve a fragmentar
  fragmentado en 3 fragmentos para el enlace con MTU 50:
    127.0.0.1;8885;18;501;0;30;1;...
    127.0.0.1;8885;18;501;30;30;1;...
    127.0.0.1;8885;18;501;60;20;1;...
\end{lstlisting}

\begin{center}
\begin{tabular}{@{}lcc@{}}
  \toprule
  Origen $\rightarrow$ destino & Fragmentos recibidos & Re-ensamblados \\
  \midrule
  R1 $\rightarrow$ R5 & 4, 4, 4, 4 & 4/4 \\
  R5 $\rightarrow$ R1 & 3, 3, 6, 6 & 4/4 \\
  R3 $\rightarrow$ R1 & 5, 5, 5, 6 & 4/4 \\
  R4 $\rightarrow$ R1 & 3, 3, 3, 3 & 4/4 \\
  R2 $\rightarrow$ R5 & 5, 4, 5, 4 & 4/4 \\
  \bottomrule
\end{tabular}
\end{center}

\textbf{Observaciones.}
\begin{itemize}[nosep]
  \item Los 20 envíos se fragmentaron y re-ensamblaron correctamente.
  \item El resultado depende del origen, porque cambian los MTU del camino. Por ejemplo,
        desde R5 el primer enlace puede tener MTU 140 (2 trozos de 120 y 10 bytes) o
        MTU 55 (trozos de 35 bytes), y por eso al destino llegaron 3 o 6 fragmentos.
  \item Sí se re-fragmentan fragmentos: el de 100 bytes que sale de R1 se dividió en 3 al
        entrar al enlace R2--R3 (MTU 50), y un trozo de 35 bytes de mensaje se dividió en
        30 y 5 bytes. Los offsets siguen siendo relativos al datagrama original.
  \item Cada fragmento pasa por el round-robin por separado, así que los fragmentos de un
        mismo paquete toman caminos distintos (entre 3 y 9 saltos) y llegan desordenados.
  \item Cada fragmento lleva su propio header de 20 bytes. En un enlace con MTU 50, el
        paquete de 150 bytes se convierte en 5 fragmentos que suman 230 bytes.
\end{itemize}

\section{Limitaciones}

\begin{itemize}[nosep]
  \item Si se pierde un fragmento, el resto de los fragmentos de su ID queda guardado
        indefinidamente en el diccionario, porque no hay un tiempo límite para
        re-ensamblar.
  \item Si el MTU es menor o igual a 20 bytes no queda espacio para datos, y el router
        descarta el paquete con un aviso.
\end{itemize}

\section*{Declaración de uso de IA}

En este informe se utilizó el modelo de IA Claude exclusivamente como apoyo en la
revisión del código, recomendaciones de diseño, la redacción del informe, la creación de
la tabla y la elaboración del diagrama en \LaTeX{}. Las decisiones de
diseño, la implementación del código, la ejecución de los experimentos y la obtención de
los datos fueron realizadas por el autor. Todo el contenido generado con apoyo de la
herramienta fue revisado y validado, y su responsabilidad es del autor.

\end{document}
